\section{Research Questions and Study Design}
\label{sec:study_design}
\label{sec:experiment}

The study asks one question: when is sample size a good enough proxy for
measured cost, and when does that shortcut change the substantive
conclusion? We organize the evaluation around three research questions.

%\begin{itemize}
    %\item 
    \textbf{RQ1. To what extent does sample size reflect measured
    acquisition cost?} We test whether same-size samples spend comparable
    measured budgets, whether larger samples reliably cost more, and whether a
    fixed fraction of configurations corresponds to the same fraction of
    measured cost.
    
    % \item 
    \textbf{RQ2. Does replacing sample size by measured cost
    change sampling-strategy conclusions?} We compare the top accuracy--cost
    choices and the Pareto-efficient choice sets under count-based cost and
    measured cost.
    
    % \item 
    \textbf{RQ3. Under which conditions is sample size a safe or unsafe
    cost proxy?} We analyze how the answer depends on cost heterogeneity,
    strategy bias, target--cost asymmetry, workload/input, release, and dataset
    family.
% \end{itemize}

\paragraph*{Operational tests.}
RQ1 first checks \emph{cost heterogeneity}. We use the coefficient of variation
(CV), the standard deviation divided by the mean, because it is unit-free: a
CV of $0.20$ means the standard deviation is $20\%$ of the mean whether the
resource is seconds, joules, or bytes, while a CV near $1$ means the spread is as
large as the average cost. CV is sensitive to broad dispersion, so we pair it
with $p_{90}/p_{10}$, the ratio between the $90$th and $10$th percentile costs.
This second statistic is robust and readable: $p_{90}/p_{10}=1.5$ means a
costly but non-extreme configuration costs $50\%$ more than a cheap but
non-extreme one, while a value near $47$ means the upper tail is about $47\times$
costlier. Fixed-$n$ strategy-effect tests then ask the operational question:
when two strategies both choose, say, $64$ configurations, do they nevertheless
visit systematically cheaper or costlier regions?

RQ1 also checks \emph{budget calibration}. Spearman $\rho$ between $n$ and
$\TrainCost$ asks only whether larger samples tend to cost more. The ratio
$\Ratio$ from \Cref{eq:ratio} asks the stricter budget question. If a system has
$1{,}000$ configurations and a sample measures $100$, then $\mathit{SizePct}=0.10$.
If those $100$ configurations spend $8\%$ of the full measured campaign, then
$\Ratio=0.08/0.10=0.8$: the sample is cheaper than its count suggests. If it
spends $15\%$, then $\Ratio=1.5$: the sample is costlier than its count
suggests.

RQ2 compares sampling conclusions under count and measured cost through three
complementary lenses, in decreasing order of how much they drive our claims.
\emph{(i)~Cost-controlled sampler comparison (primary).} The replication studies
we revisit pick the best sampler with a Kruskal--Wallis (KW) omnibus over each
sampler's per-seed error distribution and a Dunn post-hoc (Holm-adjusted, as in
the source studies) that reports the \emph{top group} of samplers not
significantly worse than the best
\cite{kruskal1952,kaltenecker2019,pereira2020}. We rerun that exact test and
change one thing: we hold the measured cost fixed instead of the
sample size. For each context we form a budget grid from the measured costs of
the predefined sample sizes. At a budget $b$ every sampler enters at the
\emph{largest predefined sample size whose median measured cost is at most~$b$},
the most data it can afford. This operating point is chosen without consulting
error, so no point is selected by test accuracy. We use it as a retrospective,
cost-indexed evaluation device, relying on observed or pilot-estimated sampler
costs rather than on test error. A sampler with no affordable size at~$b$ is excluded from
that comparison. We report whether the KW top group (the set of statistically
recommendable samplers) changes between the size- and cost-controlled views, and
how often the source study's own recommended sampler is \emph{evicted} from it.
Selecting instead the cheapest-error affordable size (a retrospective upper
envelope applied symmetrically to all samplers) leaves the conclusion unchanged
and is reported as a sensitivity check.
\emph{(ii)~Pareto choice set.} A combination is efficient if no other has both
lower error and lower measured cost. We test whether the measured-cost efficient
set differs from the count-cost one (Jaccard $<1$) and, since a single
observation could be noise, whether that difference reproduces under a
test-set-and-run bootstrap ($p_{\text{differ}}\ge0.95$). \emph{(iii)~Top-$1$
flip.} The single best strategy--budget combination under the objective $J$ may
change between count and measured cost. Because this lens considers only the
top-ranked combination, it is sensitive to small margins between near-tied
combinations, so we report it with a bootstrap robustness check and treat it as
secondary evidence. RQ3 interprets the same
statistics across contexts. We report $p$-values where they are meaningful, but
always pair them with practical thresholds: CV$\ge0.20$ or
$p_{90}/p_{10}\ge1.5$ for visible cost spread,
$|\widetilde{\Ratio}-1|\ge0.05$ for budget mismatch, and Jaccard $<1$ for a
changed efficient choice set. These thresholds are a-priori reporting
conventions fixed before analysis rather than tuned cutoffs. They only flag
where a difference is large enough to be worth reporting, and no conclusion
depends on the exact value. Full per-scope tests, $p$-values, and effect
sizes are in the supplementary artifact.

\subsection{Evidence Families}
\label{subsec:families}

We keep six evidence families separate because they stress different validity
conditions. Pooling them would give a larger number but a weaker explanation of
when the sample-size proxy is reliable. The two replication families keep the
link with published configurable-system sampling studies: one revisits seven
Kaltenecker-style systems, and the other revisits x264 over $17$ inputs and two
measured properties. The four extension families then change one condition at a
time: the cost dimension (performance versus energy), the workload or input, an
external hyperparameter-optimization corpus, and a large Linux storage-cost
setting.

The \StudyReplicationICSE family is the most conservative test. It reuses the
published strategy panel on seven Kaltenecker \etal\ systems
\cite{kaltenecker2019} ($432$--$68{,}640$ valid configurations and $11$--$54$
options each, one performance-like target also charged as cost): random,
solver-based, randomized solver-based, $t$-wise, distance-based, and diversified
distance-based sampling. It changes only the cost axis from sample count to
measured performance cost. If a count-based conclusion changes there, the change
is not caused by inventing a new sampler or a new benchmark. The
\StudyReplicationICPE family uses the same logic on x264 \cite{pereira2020}, but
adds context: the same $1{,}152$ configurations ($16$ binary options) are
observed on $17$ videos and two properties (encoding time and output size, in
four cost--target scenarios), so we can ask whether the cost proxy is safer when
target and acquisition cost are the same than when they differ.

The extension families deliberately broaden the construct, and the first two
reuse published ICSE 2023 measurement datasets rather than a sampling protocol.
We draw our own samples on their measured configurations, so any cost effect is a
property of the data and not of a re-run. \StudyEnergy asks whether the issue is
only about runtime or also appears when performance and energy can each be target
or cost (the Twins-or-False-Friends corpus \cite{weber2023twins}, fourteen systems,
$180$--$86{,}058$ configurations each, four scenarios per system). \StudyWorkload
keeps systems fixed but changes workloads, testing whether a cost-proxy conclusion
is a property of the system or of the system--workload pair (the M\"uhlbauer
\etal\ corpus \cite{muhlbauer2023workload}, nine systems, $191$--$6{,}764$
configurations, $93$ system--workload contexts, $25{,}258$ measurements, with cost
measured as wall time, or oriented throughput for H2). \StudyLCBench moves outside
configurable-software performance tables, drawing on the LCBench learning-curve
benchmark \cite{zimmer2021autopytorch}, and checks whether real runtime learning
curves in hyperparameter optimization show the same count-versus-cost gap ($35$
OpenML datasets, $2{,}000$ configurations and $7$ hyperparameters each, accuracy
as target and runtime as cost). \StudyLinux stresses scale and strategy behavior
on the TuxKConfig kernel-size datasets \cite{martin2022linux,borges2025linux}:
random, balanced, coverage, uncertainty, and diversity sample from large
randconfig pools ($19{,}952$--$92{,}562$ configurations, $494$--$497$ options
across six releases v4.15--v5.08), and the cost is the uncompressed
\texttt{vmlinux} image size, a resource-budget proxy rather than measurement time.
\subsection{Protocol}
\label{subsec:protocol}

For each available tuple of system, workload/input, target, strategy, training
size, and run, we replay or generate the training sample, compute
$\TrainCost$, train the predictor on configuration options only, and evaluate
MRE on held-out configurations. Stochastic strategies are repeated over
seeds, about $100$ runs per cell for the replication archives and $30$ for the
generated families. We exclude measured performance, energy, size, and runtime
columns from the feature matrix to avoid label leakage, especially when the
target and acquisition cost are different measured quantities.

\paragraph*{Strategy panels and diagnostic bounds.}
For the two replication families we reuse the published strategy panels and their
measured sampling traces directly: the same strategies and the same measured
configurations as Kaltenecker \etal\ \cite{kaltenecker2019} and Pereira \etal\
\cite{pereira2020}, changing only the cost axis. The generated families also add \texttt{pilotCost}, a deliberately simple
cost-aware sampler, as an existence proof rather than a tuned method. It measures
a random pilot ($20\%$ of the target size, at least four configurations, whose cost
is charged in $\TrainCost$), fits a ridge model of log cost on the configuration
options, and predicts every other configuration's cost. It then fills the sample
from the configurations predicted cheapest, by diversity sampling, so the result
stays inexpensive while spreading across the option space. Because the pilot is
charged and it uses no cost information the other samplers lack, its advantage is
a fair one. Two diagnostic bounds, \texttt{cheapOracle} and \texttt{expensiveOracle},
instead look at the true costs to pick the cheapest and the most expensive
configurations, bracketing what perfect cost-awareness could achieve. They are not
deployable and are never reported as recommendations.
